% from ega0-16-fr.tex
\oldpage[0\textsubscript{IV-1}]{27}

Enfin, pour démontrer (iii), on note que l'on a deux suites exactes
\[
  0\to{}_xM\to M\xrightarrow{x}xM\to0,\qquad
  0\to xM\to M\to M/xM\to0,
\]
et l'assertion découle aussitôt du lemme (16.2.2.2).
\end{proof}

La démonstration de (16.2.3) se fait en trois étapes~:

\begin{env}[16.2.3.2]
\label{0.16.2.3.2-fr}
On a $\dim(M)\leq d(M)$.
\end{env}

Raisonnons par récurrence sur $d(M)$, le cas $d(M)=0$ étant trivial.
Supposons donc $d(M)\geq1$, et soit $\mathfrak{p}_0\in\operatorname{Supp}(M)$
tel que $\dim(A/\mathfrak{p}_0)=\dim(M)$. Cela entraîne que
$\mathfrak{p}_0$ est un élément \emph{minimal} de $\operatorname{Supp}(M)$,
donc il est associé à $M$ (Bourbaki, \emph{Alg.~comm.}, chap.~IV, \S~1,
no~4, th.~2) et $M$ contient par suite un sous-module $N$ isomorphe à
$A/\mathfrak{p}_0$ (\emph{loc.~cit.}, \S~1, no~1)~; comme $d(M)\geq d(N)$
par (16.2.2.2), on est ramené à prouver que $\dim(N)\leq d(N)$. Soit donc
$\mathfrak{p}_0\subset\mathfrak{p}_1\subset\cdots\subset\mathfrak{p}_n$ une
chaîne d'idéaux premiers distincts dans $A$, et montrons que $n\leq d(N)$.
C'est évident si $n=0$. Sinon, il existe
$x\in\mathfrak{p}_1\cap\mathfrak{r}$ tel que $x\notin\mathfrak{p}_0$, car la
relation $\mathfrak{p}_0\supset\mathfrak{p}_1\cap\mathfrak{r}$ entraînerait
$\mathfrak{p}_0\supset\mathfrak{r}$ puisque $\mathfrak{p}_0$ ne contient pas
$\mathfrak{p}_1$, donc $\mathfrak{p}_0$ serait maximal, ce qui est absurde.
On a $\mathfrak{p}_i\in\operatorname{Supp}(N/xN)$, pour $i\geq1$, donc
$n-1\leq\dim(N/xN)$ (16.1.7). Comme ${}_xN=0$ en vertu du choix de $x$, on a
$d(N/xN)\leq d(N)-1$ par (16.2.3.1, (iii)). L'hypothèse de récurrence implique
alors $d(N/xN)=\dim(N/xN)\geq n-1$, donc $n-1\leq d(N)-1$.

\begin{env}[16.2.3.3]
\label{0.16.2.3.3-fr}
On a $d(M)\leq s(M)$.
\end{env}

Soit en effet $(x_i)_{1\leq i\leq n}$ une famille d'éléments de
$\mathfrak{r}$ telle que, si l'on pose
$\mathfrak{a}=x_1A+\cdots+x_nA$, $M/\mathfrak{a}M$ soit de longueur finie.
Les idéaux associés à $M/\mathfrak{a}M$ sont alors maximaux (Bourbaki,
\emph{Alg.~comm.}, chap.~IV, \S~2, no~5, prop.~7) et comme ce sont les seuls
idéaux premiers contenant
$\mathfrak{q}=\mathfrak{a}+(\mathfrak{r}\cap\operatorname{Ann}(M))$,
$A/\mathfrak{q}$ est artinien (\emph{loc.~cit.}, prop.~9) et $\mathfrak{q}$
est un idéal de définition de $A$. Mais on a
$\mathfrak{q}^mM/\mathfrak{q}^{m+1}M=
\mathfrak{a}^mM/\mathfrak{a}^{m+1}M$, et si $M$ est engendré par $r$ éléments
$z_j$, $\mathfrak{a}^mM/\mathfrak{a}^{m+1}M$ est un $(A/\mathfrak{q})$-module
engendré par les images canoniques des
$x_1^{\nu_1}x_2^{\nu_2}\cdots x_n^{\nu_n}z_j$ où
$\sum_i\nu_i=m$, donc par $r\binom{m+n-1}{n-1}$ éléments~; sa longueur est
par suite $\leq ar\binom{m+n-1}{n-1}$, où $a$ est une constante, et on en
déduit aussitôt que le degré de $P_{\mathfrak{q}}(M,n)$ est $\leq n$~; d'où
$d(M)\leq s(M)$.

\begin{env}[16.2.3.4]
\label{0.16.2.3.4-fr}
On a $s(M)\leq\dim(M)$.
\end{env}

Notons d'abord que $n=\dim(M)$ est fini par (16.2.3.2). Raisonnons par
récurrence sur $n$~; la proposition est évidente pour $n=0$, puisque $M$ est
alors de longueur finie (16.1.10). Supposons $n\geq1$, et soient
$\mathfrak{p}_i$ $(1\leq i\leq m)$ les idéaux premiers de
$\operatorname{Supp}(M)$ tels que $\dim(A/\mathfrak{p}_i)=n$~; la définition
des $\mathfrak{p}_i$ montre que ce sont des éléments minimaux de
$\operatorname{Supp}(M)$, donc en nombre fini (Bourbaki, \emph{Alg.~comm.},
chap.~IV, \S~1, no~4, th.~2). Comme $n\geq1$, les $\mathfrak{p}_i$ ne sont pas
maximaux, et il existe donc $x\in\mathfrak{r}$ tel que
$x\notin\mathfrak{p}_i$ pour tout $i$ (Bourbaki, \emph{Alg.~comm.}, chap.~II,
\S~1, no~1, prop.~2). En vertu de (16.2.3.1), on a
$s(M)\leq s(M/xM)+1$ et $\dim(M)\geq\dim(M/xM)+1$~; l'hypothèse de récurrence
entraîne que $s(M/xM)\leq\dim(M/xM)$, d'où la conclusion.

\begin{corollary}[16.2.4]
\label{0.16.2.4-fr}
Sous les hypothèses de (16.2.3), on a
\begin{equation}
\label{0.16.2.4.1-fr}
\tag{16.2.4.1}
  \dim_{\widehat{A}}(\widehat{M})=\dim_A(M).
\end{equation}
\end{corollary}

